\documentclass{article}
\usepackage{graphicx}% Required for inserting images
\usepackage{float}
\usepackage{amsmath}
\usepackage{booktabs}
\usepackage{hyperref}

\begin{document}
\begin{titlepage}
    \centering

    \vspace{2cm}

    {\Huge\bfseries Relazione finale progetto\par}

    \vspace{0.5cm}

    {\Large Analisi del rischio di default \par}

    \vspace{2cm}
    
    {\Large 
    \href{https://github.com/caterinaghiani/progetto_datascience}{\textbf{Repository GitHub}}
    \par}

    \vspace{1.5cm}

    {\large
    Caterina Ghiani - 1181102
    }

    \vfill

    {\large Anno Accademico 2025/2026 \par}
\end{titlepage}


\thispagestyle{empty}
\newpage
\tableofcontents
\thispagestyle{empty}
\newpage

\setcounter{page}{1}
\section{Introduzione}
\textbf{Domanda di ricerca} -- \`E possibile prevedere se un cliente andr\`a in default sul pagamento del mese successivo ($Y=1$) analizzando il suo profilo demografico, il limite di credito concesso, lo storico dei pagamenti e l'ammontare degli estratti conto passati?
\\
Per le banche, la valutazione del rischio di credito dei propri clienti \`e fondamentale: essa consente di identificare in anticipo i clienti a rischio di insolvenza, riducendo di conseguenza le perdite. Individuare pattern di comportamento ricorrenti permette inoltre di delineare il profilo dei soggetti potenzialmente meno affidabili.
\\
La relazione descrive l'intero processo di analisi condotto, articolato in tre fasi: 
\begin{itemize}
    \item analisi esplorativa univariata e bivariata 
    \item modellazione
    \item confronto della performance dei modelli
\end{itemize}

\section{Dataset}

Il dataset proviene da una banca, la quale ha fornito informazioni riguardo i titolari di carte di credito per prevedere se il cliente e titolare della carta, pagher\`a o meno il suo prossimo accredito. Per poter condurre tale previsione, il dataset contiene informazioni sulle caratteristiche anagrafiche dei clienti quali ad esempio, genere e et\`a, ma anche il limite di credito, lo storico dei pagamenti e i comportamenti assunti in precedenza. 
\\
\subsection{Descrizione delle variabili}
Il dataset ha 30.000 osservazioni: si riferisce a ben 30.000 individui e ha 24 features. Non sono stati trovati dati mancanti.
Rispettivamente le variabili contengono:
\begin{itemize}
    \item \textbf{LIMIT\_BAL} (numerica discreta) -- \`e l'importo del credito concesso (in dollari) e include sia il credito al consumo individuale sia quello familiare (supplementare)
    \item \textbf{SEX} (booleana) -- il genere del titolare, che assume valore pari a 0 per i maschi e 1 per le femmine
    \item \textbf{EDUCATION} (categoriale ordinale) -- l'istruzione del titolare: 1 = specialistica/master/dottorato, 2 = universit\`a, 3 = scuola superiore e 4 = altro
    \item \textbf{MARRIAGE} (categoriale nominale) -- lo stato civile del titolare: 1 = sposato/a, 2 = single e 3 = altro
    \item \textbf{AGE} (numerica discreta) -- l'et\`a del titolare espressa in anni
    \item \textbf{PAY\_X} (categoriale ordinale) -- \`e lo stato del rimborso a settembre (\texttt{PAY\_1}) fino a \texttt{PAY\_6}, che \`e lo stato del rimborso ad aprile; la scala di misurazione dello stato del rimborso \`e: $-1$ e $0$ pagamento regolare, $1$ in ritardo di un mese, fino a $9$ in ritardo di nove o pi\`u mesi
    \item \textbf{BILL\_AMT} (numerica discreta) -- \`e l'importo dell'estratto conto a settembre (\texttt{BILL\_AMT1}), fino a \texttt{BILL\_AMT6}, che \`e l'importo dell'estratto conto ad aprile
    \item \textbf{PAY\_AMT} (numerica discreta) -- \`e l'importo del pagamento precedente, in questo caso l'importo pagato a settembre (\texttt{PAY\_AMT1}), fino a \texttt{PAY\_AMT6}, che \`e l'importo pagato ad aprile
    \item \textbf{default payment} (booleana, variabile target) -- \`e pari a 1 se il titolare ha fatto default e 0 se ha pagato
\end{itemize}

\section{Metodologia}

\textbf{Pulizia dei dati}: \`e stato necessario rinominare la variabile \textit{PAY\_0} in \textit{PAY\_1} per garantire la sequenzialit\`a temporale con i mesi successivi. Per coerenza sono stati accorpati in una classe unica $0$ ("Nessun ritardo") i valori negativi presenti nella variabile \textit{PAY\_N}, riservando i valori interi positivi ($1, 2, \dots$) ai mesi effettivi di ritardo nei pagamenti.\\
\textbf{Gestione assenza dati e outlier}: in assenza di valori nulli ($30.000$ osservazioni complete), si \`e deciso di mantenere gli outlier finanziari per preservare i casi estremi di rischio effettivo.\\
\textbf{Preparazione dei dati e Feature Engineering}:utilizzando \textit{StandardScaler}, sono state standardizate le variabili numeriche per portarle a media 0 e deviazione standard 1, per k-NN e regressione logistica. \`E stata creata la variabile \textit{UTILIZATION\_RATIO} $= (\text{BILL\_AMT1} / \text{LIMIT\_BAL}) \times 100$ per misurare la percentuale di saturazione del tetto di spesa ed eliminare la distorsione di scala dei debiti assoluti.
  

\subsection{Modelli Utilizzati}
\begin{itemize}
    \item \textbf{Regressione Logistica}: ci permette di analizzare la curva sigmoide e capire come l'aggiustamento della soglia decisionale aiuti a gestire il rischio di default.
    \item \textbf{k-Nearest Neighbors (k-NN)}: modello non parametrico che classifica i clienti in base alla vicinanza ai loro $k$ profili pi\'u simili nello spazio delle feature.
    \item \textbf{Random Forest}: modello \textit{ensemble} basato su 200 alberi decisionali. Addestrato sui dati originali (non scalati), \`e ideale per comprendere relazioni non lineari e per estrarre la \textit{feature importance}.
\end{itemize}

\subsection{Validazione e Partitioning}
\begin{itemize}
    \item \textbf{Train/Test Split}: suddivisione 80/20 del dataset ($24.000$ clienti per il training e $6.000$ per il test).
    \item \textbf{Stratificazione}: utilizzata l'opzione \textit{stratify=y} per conservare la proporzione originale dei default (approssimativamente $22.1\%$) sia nell'addestramento che nella verifica.
    \item \textbf{Prevenzione Data Leakage}: lo \textit{StandardScaler} \`e stato calibrato (\textit{fit\_transform}) solo sui dati di train ed esteso (\textit{transform}) al test set.
\end{itemize}


\section{Risultati}
\subsection{Risultati dell'analisi esplorativa}
\subsubsection{Analisi univariata}
\begin{enumerate}
    \item \textbf{Distribuzione della variabile target $Y$}, i 3/4 dei soggetti osservati \`e solvibile, mentre 1/4 \`e in default. Questo indica che il nostro dataset non \`e bilanciato rispetto alla variabile target
    \begin{table}[htbp]
    \centering
    \begin{tabular}{ccc}
    \toprule
    \textbf{} & \textbf{Conteggio (N)} & \textbf{Percentuale (\%)} \\
    \midrule
    \textbf{Solvibile (0)} & 23364 & 77.88 \\
    \textbf{Default (1)} & 6636 & 22.12 \\
    \bottomrule
    \end{tabular}
    \caption{distribuzione della variabile target $Y$}
    \label{tab:mia_tabella}
    \end{table}

    \item \textbf{Caratteristiche anagrafiche dei nostri soggetti}: il 60\% del nostro campione \`e composto da soggetti di genere femminile. Circa la met\`a delle persone ha ricevuto un'educazione universitaria e il 35\% del campione ha continuato gli studi dopo l'universit\`a. Inoltre, rispetto alla situazione sentimentale, ci troviamo davanti ad un campione in cui poco pi\`u della met\`a dei soggetti \`e single e l'altra met\`a \`e sposata (l'1\% non ha voluto dichiararlo o non si identifica in nessuna delle categorie). L'et\`a dei soggetti \`e compresa tra i $21$ e i $79$ anni, ma 3/4 dei soggetti hanno meno di 41 anni. La media \`e di 35 anni, questo vuol dire che probabilmente il valore 79 rappresenta un outlier.

    \item \textbf{Distribuzione del limite di credito accordato}: la media del credito \`e di 167.484,32 dollari, con una forte deviazione standard. La mediana \`e inferiore alla media, evidenziando una distribuzione asimmetrica a destra. Difatti, pochi clienti hanno crediti molto elevati e lo possiamo notare anche dal massimo del credito concesso, che \`e pari a 1 milione.

    \item \textbf{Analisi dello storico dei pagamenti}: calcolando la distribuzione percentuale dei mesi di ritardo ricostruendo lo storico dei pagamenti, individuiamo che circa l'80\% dei soggetti \`e in pari con i pagamenti della propria carta di credito, con un picco del 90\% nel mese di aprile. Si nota anche che i soggetti in ritardo sono solitamente in ritardo di 2 mesi nei pagamenti da effettuare. 

    \item \textbf{Importo delle bollette mensili}: la media dell'importo delle bollette aumenta man mano che si progredisce nell'anno e in tutti i mesi la media \`e pi\`u del doppio della mediana. Il 50\% dei clienti ha bollette contenute e approssimativamente degli stessi importi, mentre pochi clienti accumulano saldi estremamente elevati. I valori negativi non indicano un errore ma implicano un credito a favore del cliente.
    
\end{enumerate}

\subsubsection{Analisi bivariata}
\begin{enumerate}
    \item \textbf{L'et\`a differisce in modo significativo tra clienti in default e clienti solvibili?}\\ Si rileva una mediana pari a circa $34$ anni, con un range interquartile compreso tra 28 e 41 anni. La presenza di valori anomali statistici sopra i 60 anni rispecchia quanto visto nell'analisi univariata in cui pochi soggetti hanno un'et\`a avanzata ma non costituiscono un errore, piuttosto rappresentano una situazione realistica. 
    La quasi totale somiglianza tra i due boxplot, suggerisce che l'et\`a , non \`e un fattore che influisce sulla variabile target.
    \begin{figure}[H]
    \centering
    \includegraphics[width=1.1\linewidth]{boxplot_age_target.png}
    \caption{Distribuzione dell'et\`a rispetto alla variabile target $Y$}
    \label{fig:placeholder1}
    \end{figure}

    \item \textbf{Un elevato tasso d'istruzione si associa ad una minore probabilit\`a di default?}\\
    Al diminuire del grado del titolo di studio aumenta il tasso di default: ad un'istruzione pi\`u elevata si associa una maggiore probabilit\`a di solvibilit\`a supponendo che sia legata ad un reddito medio pi\`u stabile.
    I clienti codificati come \textit{``Altro''} includono anche le anomalie che sono state accorpate e registrano un tasso di default estremamente basso (7.1\%). Questo potrebbe essere dovuto al fatto che corrispondano a categorie particolari, probabilmente a categorie di soggetti sui quali la stessa banca fa uno screening pi\`u approfondito per verificarne l'affidabilit\`a.
    \begin{figure}[H]
    \centering
    \includegraphics[width=1.1\linewidth]{education_stacked_breakdown.png}
    \caption{Ripartizione percentuale di solvibili e insolventi per titolo di studio}
    \label{fig:placeholder2}
    \end{figure}

    \item \textbf{Esiste una differenza in base al genere che rende i soggetti pi\`u o meno probabili di essere in default?}\\
    Calcolando il tasso medio di default per genere possiamo confrontare le percentuali normalizzate: le donne hanno un tasso di insolvenza inferiore rispetto agli uomini (20.85\% donne e 24.17\% uomini) quindi le donne sono pi\`u affidabili degli uomini. A causa della sproporzione nel portafoglio clienti (le donne sono il 60\%), in valore assoluto le donne superano gli uomini nei casi di default. 
    \begin{table}[H]
    \centering
    \begin{tabular}{lccccc}
    \toprule
    \textbf{Genere} & \textbf{Solvibili (N)} & \textbf{Solvibili (\%)} & \textbf{Default (N)} & \textbf{Default (\%)}  \\
    \midrule
    Maschio (0) & 9015  & 75.83 & 2873 & 24.17  \\
    Femmina (1) & 14349 & 79.22 & 3763 & 20.78  \\
    \bottomrule
    \end{tabular}
    \caption{Distribuzione del default per genere (valori assoluti e percentuali)}
    \label{tab:default_genere}
    \end{table}

    \item \textbf{Un soggetto sposato \`e pi\`u probabile che sia solvibile?} \\
    le categorie ``sposato/a" e ``single" sono approssimativamente bilanciate avendo rispettivamente $13.659$ e $15.964$ soggetti osservati. Questo bilanciamento lo notiamo anche nella distribuzione delle due categorie rispetto alla variabile target. Possiamo dire che tra la variabile target e lo stato civile non vi \`e una forte correlazione.

    \item \textbf{Un limite di credito concesso pi\`u basso segnala un soggetto pi\`u a rischio di default?}\\
    Il grafico mostra una tendenza decrescente: all'aumentare del credito accordato, la percentuale di clienti in default si riduce in modo costante.
    \begin{itemize}
        \item \textit{Fascia ad alto rischio ($<$ 50k)} in cui i clienti hanno un tasso di default pi\`u alto di tutti: quasi 1 cliente su 3 in questa fascia va in default;
        \item \textit{Fascia intermedia (50k - 300k)} il rischio di default diminuisce progressivamente dal 25.8\% al 16.1\%;
        \item \textit{Fascia a basso rischio ($>$ 300k)} questi clienti a cui la banca d\`a molta fiducia, effettivamente hanno una percentuale bassa di default, di 2.4 volte pi\`u affidabili della prima fascia. Le banche tendenzialmente verificano in modo approfondito l'affidabilit\`a di clienti prima di concedere limiti di credito cos\'i alti.
    \end{itemize}
    
    \begin{figure}[H]
    \centering
    \includegraphics[width=1.1\linewidth]{limit_bal_fasce_numerosita.png}
    \caption{Tasso di default per fascia di limite di credito}
    \label{fig:placeholder3}
    \end{figure}

    \item \textbf{Un elevato rapporto tra spesa mensile e credito concesso aumenta la probabilit\`a di default? } \\
    Dal rapporto dell'ultimo estratto conto (BILL\_AMT1) e il limite di credito (LIMIT\_BAL) guardiamo il debito in rapporto alla capacit\`a di ripagarlo. I soggetti finanziariamente pi\`u vulnerabili sono coloro che lavorando costantemente a ridosso del tetto massimo di spesa, non dispongono di alcun cuscinetto di liquidit\`a residua. Di conseguenza, di fronte a un piccolo imprevisto o a un'emergenza, la loro probabilit\`a di andare in default aumenta notevolmente.
    
    \begin{figure}[H]
    \centering
    \includegraphics[width=1.1\linewidth]{credit_utilization_rates.png}
    \caption{Tasso di default per fascia di limite di utilizzo del credito}
    \label{fig:placeholder4}
    \end{figure}
    
    In questo grafico osserviamo:
    \begin{itemize}
        \item Da 1\% a 30\% \`e la fascia pi\`u stabile, perch\'e \`e numerosa e con il tasso di default pi\`u basso in assoluto
        \item dal 71\% in su invece abbiamo un tasso di default pi\`u importante, per cui circa 9 mila soggetti sono ad alto rischio e vivono una situazione di stress finanziario
        \item anche coloro che sono inattivi, sono per\`o un grande rischio perch\'e sono clienti dormienti.
    \end{itemize}
\end{enumerate}

\subsection{Tecniche predittive}
\subsubsection{Regressione logistica}
    Osservando la distribuzione della variabile target rispetto alla variabile\\ \textit{UTILIZATION\_RATIO} si nota visivamente una maggiore concentrazione di casi di default ($Y$=1) all'aumentare del livello di saturazione del tetto di credito, suggerendo una relazione positiva tra l'elevato utilizzo del credito e l'insolvenza. 
    Il coefficiente \`e infatti positivo ($0.005$) sottolineando la relazione tra saturazione del credito e rischio di default. L'intercetta \`e invece negativa ($-1.48$) questo perch\'e quando un cliente non utilizza il suo credito ha comunque probabilit\`a di default, un dato che si osserva gi\`a nell'analisi bivariata.\\

    Utilizzando tre soglie differenti: 0.5, 0.3, 0.2 identifichiamo la soglia prudenziale oltre la quale il modello inizia a classificare come a rischio di default i soggetti. In questo caso la soglia prudenziale \`e di 0.2, cio\`e tutti i clienti che superano l'utilizzo del credito del 30\% sono classificati come a rischio. Unicamente il cliente con un utilizzo del 10\% viene considerato completamente solvibile. 
        \begin{table}[H]
        \centering
        \begin{tabular}{cccccc}
        \toprule
        \textbf{Indice} & \textbf{UT.\_RATIO} & \textbf{prob. default} & \textbf{soglia\_05} & \textbf{soglia\_03} & \textbf{soglia\_02} \\
        \midrule
        0 & 10  & 0.1933 & 0 & 0 & 0 \\
        1 & 30  & 0.2091 & 0 & 0 & 1 \\
        2 & 70  & 0.2433 & 0 & 0 & 1 \\
        3 & 95  & 0.2666 & 0 & 0 & 1 \\
        4 & 120 & 0.2912 & 0 & 0 & 1 \\
        \bottomrule
        \end{tabular}
        \caption{Predizioni di default al variare della soglia di classificazione}
        \label{tab:predizioni_soglia}
        \end{table}

\subsubsection{k-Nearest Neighbour}
    Con questo modello \`e necessaria la standardizzazione perch\'e basandosi sulle distanze, variabili espresse in valuta reale come $LIMIT\_BAL$, avrebbero dominato la misura di distanza rispetto a features come $SEX$ o $MARRIAGE$. \\
    Per individuare la configurazione migliore del k-NN, ne abbiamo testati diversi. 
    \begin{itemize}
        \item k troppo bassi (\textit{k = 1, 3}) il modello \`e troppo sensibile al rumore dei dati tenendo un'accuracy bassa (con k = 1, accuracy = 73\%)
        \item scelta ottimale (\textit{k = 21}) allargando il raggio di K di molto, le predizioni sono pi\`u stabili e l'accuracy raggiunge un valore di 81\%
        \item k troppo alti (\textit{k = 51}) oltre i 21 le prestazioni iniziano a calare, poich\'e il modello tende ad appiattirsi sulla classe maggioritaria (clienti solvibili)
    \end{itemize}
        \begin{figure}[H]
        \centering
        \includegraphics[width=1.1\linewidth]{k_NN.png}
        \caption{Effetto parametro k}
        \label{fig:placeholder5}
        \end{figure}

\subsubsection{Random Forest}
    A differenza dei modelli precedenti, l'addestramento \`e  stato eseguito sui dati originali non scalati, sfruttando la naturale insensibilit\`a degli alberi alle trasformazioni monotone delle features. \`E stata addestrata una random forest composta da 200 alberi decisionali.\\
    La classifica dell'importanza delle variabili, individua come predittori chiave: lo storico dei pagamenti recenti $PAY\_1$, l'et\`a $AGE$ e l'indice di utilizzo del credito $UTILIZATION\_RATIO$. Utilizzando una variabile $random\_noise$, che si introduce come feature casuale avente nessuna correlazione con il nostro dataset, ci nota come solo $PAY\_1$ \`e  effettivamente importante come predittore di rischio, siccome come tutto ci\`o che sta al di sotto della variabile casuale non sia cos\'i rilevante. 

    \begin{figure}[H]
    \centering
    \includegraphics[width=1.1\linewidth]{feauture_importance2.png}
    \caption{Feature importance con variabile casuale}
    \label{fig:placeholder6}
    \end{figure}

\subsection{Performance dei modelli}
    Tutti e tre i modelli mostrano un'accuracy quasi identica. Questo dato conferma come l'accuracy sia una metrica ingannevole quando si \`e in presenza di un dataset sbilanciato.\\
    La regressione logistica ottiene la precision pi\`u alta, in quanto identifica il default correttamente 2 volte su 3. Ha tuttavia una recall pi\`u bassa (0.33) lasciandosi quindi sfuggire ben il 67\% dei clienti a rischio. \`E meglio avere pi\`u falsi positivi piuttosto che sottostimare il rischio di default.\\
    Random forest mostra la capacit\`a di intercettazione migliore con una recall del 0.37 ma ha una precision leggermente inferiore agli altri modelli.\\
    Il modello k-NN si colloca esattamente a met\`a tra i due con una precision del 0.65 e una recall di 0.35.\\
        \begin{table}[H]
        \centering
        \begin{tabular}{lccc}
        \toprule
        \textbf{Metrica} & \textbf{Regressione Logistica} & \textbf{k-NN (k=21)} & \textbf{Random Forest} \\
        \midrule
        Accuracy   & 0.82 & 0.81 & 0.81 \\
        Precision  & 0.67 & 0.65 & 0.63 \\
        Recall     & 0.33 & 0.35 & 0.37 \\
        F1-score   & 0.44 & 0.45 & 0.46 \\
        \bottomrule
        \end{tabular}
        \caption{Confronto delle metriche di performance tra i modelli predittivi}
        \label{tab:confronto_modelli}
        \end{table}



\section{Discussione}
    L'analisi evidenzia come lo storico dei pagamenti $PAY\_1$ e l'indice di saturazione del credito $UTILIZATION\_RATIO$ siano predittori efficaci per stimare il rischio di default. L'introduzione di quest'ultimo indice riassuntivo \`e  risultata una scelta vincente per descrivere il comportamento dei soggetti osservati.\\
    Il principale vincolo metodologico \`e  rappresentato dallo sbilanciamento del dataset (circa il 22\% di casi default). In questo contesto, l'accuracy risulta una metrica fuorviante, mentre la priorit\`a per la banca \`e  massimizzare la recall per ridurre i falsi negativi. L'adozione della random forest e il ricalibro della soglia decisionale si confermano strategie necessarie per intercettare una quota maggiore di clienti a rischio.\\
    Infine, il processo ha evidenziato l'importanza della pulizia dei dati per correggere le incongruenze nella documentazione originale. Nonostante la natura semplificata dello studio, i risultati dimostrano l'efficacia del machine learning nel supportare decisioni operative di gestione del rischio del credito. 

\section{Conclusione}
    In sintesi, questo lavoro ha analizzato i dati demografici e finanziari di 30.000 clienti bancari, applicando tre algoritmi per la classificazione per prevedere il rischio di default sul pagamento del mese successivo.\\
    Alla domanda di ricerca che ci siamo posti inizialmente troviamo un riscontro positivo. \`E possibile prevedere il default analizzando il profilo del cliente. Risultano chiave in particolare lo storico dei pagamenti e l'indice di utilizzo del credito. Tra i modelli, il random forest \`e  risultato la scelta migliore per prevedere il maggior numero di clienti a rischio, limitando i falsi negativi.\\
    Per migliorare ulteriormente le prestazioni in futuro, si potrebbe utilizzare un campione pi\`u bilanciato e aggiungere informazioni rispetto alla situazione finanziaria del cliente e sperimentare con altri algoritmi.

\bibliographystyle{plain}
\nocite{*}
\bibliography{references}

\end{document}
